% Modelo de Relatório de Pesquisa - UFAPE.
% Baseado no Template para Relatório de Pesquisa.docx
% (docs-originais) e padronizado no estilo dos modelos de
% Monografia e ESO (mesma classe ufape, mesmas margens, fontes e títulos).
% No Overleaf/Prism, abra o main.tex e compile.
% Opções: arial (fonte Arial), notas (instruções do Word), frenteeverso (impressão dupla face).
\documentclass[monografia]{ufape}

% Dados
\autor{Nome do Estudante}
\titulo{TÍTULO DO PLANO DE TRABALHO}
\tituloprojeto{TÍTULO DO PROJETO DE PESQUISA: XXXXX}
\programa{(NOME DO PROGRAMA)}
\areapesquisa{XXXX}
\tiporelatorio{Relatório de XXX}
\orientador{Prof. Dr. Xxxxxxxxxxxxxxxxx}
\periodo{xx/xx/xxxx a xx/xx/xxxx}
\cidade{Garanhuns}
\ano{2026}

\begin{document}

% Pré-textuais
\caparelatorio
\folhaderostorelatorio

% Orientações gerais do Word. Aparecem só com a opção `notas`; na versão
% final elas somem (os campos de instrução devem ser apagados). A caixa
% aparecia duplicada no Word; aqui ela entra uma única vez.
\nota{A elaboração do relatório deve obedecer à estrutura definida neste modelo, com limite máximo de 20 páginas, contadas a partir da seção de Identificação do Projeto. O texto deve ser digitado em fonte Times New Roman, tamanho 12, com alinhamento justificado e espaçamento entre linhas de 1,5 salvo quando houver orientações específicas para campos determinados.}
\nota{Os campos com instruções ou exemplos de preenchimento devem ser integralmente apagados antes da versão final do documento. Recomenda-se atenção à clareza, coesão textual e correção gramatical ao longo de toda a produção.}
\nota{Todas as seções do relatório devem estar devidamente numeradas e identificadas com títulos em destaque, preferencialmente em negrito, a fim de garantir organização, legibilidade e fácil localização das informações pelo leitor ou avaliador.}

\sumario

% Textuais
\iniciotexto

\chapter{Introdução}
Esta seção deve apresentar o contexto e a justificativa do estudo, destacando
a relevância do tema abordado. Essa seção também pode definir o problema de
pesquisa e as hipóteses a serem testadas, estabelecendo o direcionamento do
trabalho.


\chapter{Objetivos}
A seção de objetivos geral e objetivos específicos em um relatório de pesquisa
tem como finalidade esclarecer o que se pretende alcançar com o estudo.

\section{Objetivo geral}

O objetivo geral descreve de forma ampla e concisa o propósito principal da
pesquisa.

\section{Objetivos específicos}

Os objetivos específicos detalham as metas mais concretas e mensuráveis que,
quando alcançadas, contribuem para a realização do objetivo geral.


\chapter{Fundamentação teórica}
Esta seção deve revisar os principais conceitos, teorias e estudos anteriores
relacionados ao assunto, fornecendo uma base teórica sólida que embasa a
pesquisa.


\chapter{Metodologia}
Esta seção deve descrever as etapas, métodos, técnicas e procedimentos
utilizados para alcançar os objetivos do estudo. Isso inclui o tipo de pesquisa
(qualitativa, quantitativa ou mista), a abordagem adotada, os critérios de
seleção da amostra, os instrumentos de coleta de dados (questionários,
entrevistas, experimentos, etc.), as técnicas de análise empregadas e de
tratamento dos dados. Essa seção garante a transparência do estudo e permite
que outros pesquisadores possam replicá-lo ou avaliar sua validade.

\nota{\textbf{Para o Relatório Parcial:} descreva nesta seção as atividades
metodológicas realizadas durante o período entre o início da vigência do Plano
de Trabalho e a entrega do Relatório Parcial.}


\chapter{Resultados e discussão / Desenvolvimento do trabalho}
Esta seção apresenta os principais achados do estudo e sua interpretação à luz
da fundamentação teórica. Os resultados devem ser expostos de forma clara,
utilizando tabelas, gráficos ou descrições conforme necessários. Na discussão,
esses achados são analisados e comparados com estudos anteriores, destacando
suas implicações, limitações e possíveis explicações. Essa seção deve
evidenciar se os objetivos da pesquisa foram alcançados e contribuir para o
entendimento do tema investigado.

\nota{\textbf{Para o Relatório Parcial (resultados parciais):} se a pesquisa
avançou a ponto de ter alguns resultados, apresente-os aqui. Faça uma breve
discussão sobre o que foi alcançado até o momento. Utilize tabelas, gráficos,
fotos, figuras e ilustrações que facilitem a apresentação ordenada dos
resultados.}

\section{Xxxxxx}

Abaixo seguem um exemplo de subseção e figura com suas formatações. Em seguida
apresenta-se um exemplo de tabela (figura~\ref{fig:exemplo} e
tabela~\ref{tab:exemplo}). Troque o conteúdo de exemplo pelos seus resultados.

\begin{figure}[H]
  \caption{Ponto de colheita do gladíolo no estágio R2 da Escala PhenoGlad de desenvolvimento.}
  \label{fig:exemplo}
  \centering
  \begin{minipage}{320pt}
    \centering
    \fbox{\parbox[c][5cm][c]{300pt}{\centering (Substitua esta caixa pela figura)}}\par
    \fonte[4pt]{Adaptada de xxxxx et al. (2019).}
  \end{minipage}
\end{figure}

\begin{table}[H]
  \caption{xxxxxxxxx}
  \label{tab:exemplo}
  \tabelaibge\raggedright
  \setlength{\tabcolsep}{5pt}\renewcommand{\arraystretch}{1.3}\setlength{\arrayrulewidth}{0.5pt}
  \noindent\begin{tabular*}{\textwidth}{@{\extracolsep{\fill}}|c|c|c|c|}
    \hline
    \textbf{Coluna 1} & \textbf{Coluna 2} & \textbf{Coluna 3} & \textbf{Coluna 4} \\ \hline
    xxxxx & xxxxx & xxxxx & xxxxx \\ \hline
    xxxxx & xxxxx & xxxxx & xxxxx \\ \hline
    xxxxx & xxxxx & xxxxx & xxxxx \\ \hline
    xxxxx & xxxxx & xxxxx & xxxxx \\ \hline
    xxxxx & xxxxx & xxxxx & xxxxx \\ \hline
    xxxxx & xxxxx & xxxxx & xxxxx \\ \hline
    xxxxx & xxxxx & xxxxx & xxxxx \\ \hline
  \end{tabular*}\par
  \fonte{xxxxx (ano)}
\end{table}


\chapter{Considerações parciais ou finais}
Esta seção deve sintetizar os principais achados do estudo, destacando se os
objetivos foram alcançados e quais contribuições foram feitas para a área. Além
disso, pode apontar limitações da pesquisa e sugerir direções para estudos
futuros. Nas considerações finais, ela conclui o estudo de forma definitiva,
reforçando sua relevância e impacto.

\nota{\textbf{Para o Relatório Parcial:} no caso de considerações parciais, a
seção apresenta um fechamento provisório, indicando a continuidade do trabalho.}


\chapter{Dificuldades encontradas}
Esta parte deve apontar eventuais empecilhos enfrentados durante as atividades.
Por exemplo, indisponibilidade temporária de laboratório, problemas de saúde,
adversidades no campo de pesquisa e etc (dificuldades ou limitações quanto à
geração e coleta de dados).


% Pós-textuais (a seção 8 do Word; sem número, como na monografia e no ESO)
\begin{referencias}
  \refitem{AGNOLETTI, B.Z.; PINHEIRO, P.F.; PEREIRA, L.L.; OLIVEIRA, E.C.S.;
FILGUEIRAS, P. R. Discriminação de Café Arábica Quanto à Qualidade Sensorial
Utilizando Cromatografia Gasosa e PLS-DA. \textbf{Revista Ifes Ciência},
[S. l.], v. 6, n. 3, p. 147-158, 2020.}

\refnota{Liste todas as fontes citadas ao longo do trabalho, em ordem
alfabética, alinhadas à esquerda, em espaço simples e com uma linha em branco
entre elas, conforme a última versão publicada da NBR ABNT 6023. Apague esta
nota e a referência de exemplo ao escrever.}

\end{referencias}

% As referências ocupam a posição 8; os capítulos retomam em 9 e 10.
\setcounter{chapter}{8}

\chapter{Atividades do plano de trabalho não realizadas}
Esta parte especifica as atividades planejadas no Plano de Trabalho, mas não
executadas. Uma breve justificativa deve ser explicada/proposta.


\chapter{Parecer do(a) orientador(a)}
Aqui deve ter um texto de anuência do(a) orientador(a).

\vspace{20pt}
\begin{flushright}
Garanhuns, \underline{\hspace{0.9cm}} de \underline{\hspace{2.5cm}} de 20\underline{\hspace{0.6cm}}.
\end{flushright}

\vspace{40pt}
\begin{flushright}
\begin{minipage}{7.5cm}
  \centering
  \rule{7.5cm}{0.4pt}\\
  Assinatura do(a) Orientador(a)
\end{minipage}
\end{flushright}

\vspace{30pt}
\begin{flushright}
\begin{minipage}{7.5cm}
  \centering
  \rule{7.5cm}{0.4pt}\\
  Assinatura do(a) Discente
\end{minipage}
\end{flushright}


\end{document}
